\pagelayout{wide}
% \setchapterpreamble[u]{\margintoc}
\chapter*{Introduction\,\,}
\labch{intro}
\addcontentsline{toc}{chapter}{Introduction}

% % TODO: Subsection label <05-05-23, kunzaatko> 
Advancements in precise and accurate observation techniques have always been at the forefront of scientific endeavours.
Novel approaches often lead to chains of revolutionary new discoveries in the fields that utilize these methods.

In microscopy, a physical limit that for a long time seemed unbreakable was \nameemph{Abbe's} resolution limit, due to
    diffraction induced by the light transfer through the microscope.
This bound, restricting both the lateral (\(\approx \qty{180}{\nano\meter}\)) and the axial 
    (\(\approx \qty{450}{\nano\meter}\)) axes of resolution, remained a constraint for the precision of observation in 
    microscopy for more than a century from its derivation, stifling progress in biology and countless other fields 
    which depend on measurement and observation through microscopy imaging.

Hope for better resolving power was re-established by a breakthrough at the hands of the \nameemph{Cremer} brothers in
    the mid 1970s by a revolutionary approach to microscopy observation (\acs{4PIM}) that lead to the possibility of 
    about \(4\)-fold improvement in the axial resolution.
This discovery incited a revolution, in which new methods where introduced at an unprecedented speed for this field.
As a result of this rise of attention, not long after, methods for improving the more practicable lateral resolution 
    where developed and implemented.
These methods which outperform the physical limits of \nameemph{Abbe} in at least one of the axes of resolution became 
    collectively known as \acf{SR} methods.
Most of these methods, however, relied on long acquisition times stemming from the necessity of sensing a narrow volume
    of the sample object field which deemed them unusable for applications with non-static specimen.

This pressing barrier for implementing \ac{SR} for use in \emph{in vivo} examination was overcome with a new technique 
    in the mid 1990s which was based on significantly a different principles within the optical acquisition than the
    previous methods.
Improvements that built upon the presented aforementioned \ac{SIM} method were rapidly developed and its implementation
    was widely adopted throughout the scientific community.
Further methods that employed the principles of \ac{SIM} also found there grounds in application.
These two waves of research have resulted in the creation of a plethora of approaches and implementations of \ac{SR} 
    methods.

The original \ac{SIM} technique is still very relevant in present-day and the reconstruction methods for it are
    continually improving.
An underlying concept behind \ac{SIM} is the transfer function, which is modulated by the \emph{structured illumination}
    which the \ac{SIM} method prescribes.
With respect to this, the knowledge of the correct transfer function is fundamental to reconstruction techniques of the
    \ac{SIM} method.

A transfer function is a powerful concept, which is not limited to microscopy.
It is fundamental due to the fact that it fully determines the mapping of a linear system, which it represents.
The utility of the knowledge of transfer functions is undeniable since physical systems are often linear or can be
    approximated by a linear system.

In the context of \ac{SIM}, an incorrect transfer function, may lead to artefacts and undetermined fidelity of the 
    reconstruction and so it is crucial to determine it correctly.
With this in mind, there are two major approaches to deriving the transfer function of the optical system:
    \begin{description}
        \item[modeling] Considers the context of the optical setup and the nature of the propagating light to determine
            an appropriate transfer function with the relevant aberrations of the lens system and aperture.
        \item[estimation] Resorts to the sensed signal of an acquisition and through modeling the source signal which 
            produced the result, infers an estimated transfer function.
    \end{description}

I use the estimation approach to derive the applicable transfer function of our acquisition in approximation for the 
    whole field and subsequently for the individual locations within the focal plane.
I incorporate this estimate into the reconstruction procedure to achieve better results than were previously available 
    for use in our analysis of endocytosis at the outer cell membrane.

% Their charm rises from the possibility of viewing the known with a greater closeness and opening the prospect of
%     discovering what was impossible with previous methods, enabling the recognition of new phenomena in the field.
%
% % \section{Motivation}
% \section*{The Need for Precision}
% \labsec{need-for-precision}
% \addcontentsline{toc}{section}{Need for Precision}  
% % TODO: Example with Tycho Brahe - need for greater precision <05-05-23, kunzaatko> 
%
% We can effortlessly follow the complementary relationship between significant discoveries and the conceptions of novel
%     observation approaches throughout the history of science.
% % TODO: margin note of the past most known discoveries after observational novelties <26-03-23> 
% The prominence and crucial position of the development of sensing methods, within the scientific disciplines, has not
%     diminished over time.
% % TODO: margin note on Nobel prizes in the past decade for sensing techniques <26-03-23, kunzaatko> 
% In fact the current direction of research towards being a combined effort of progressively greater groups with broader
%     specializations permits teams to appropriate greater means to fine-tuning the data acquisition procedures used. 
% % TODO: Add reference <29-03-23> 
%
% % TODO: More fluent progression to next part <29-03-23> 
% Game changing breakthroughs in sensing, however appealing, are not a matter of daily occurrence in science and with
%     this in mind scientists are expending serious effort on the improvement of techniques that are the current state of
%     the art in the field.
%
% % TODO: Figure with a timeline of significant technology of sensing discoveries. <26-03-23, kunzaatko> 
%
% Although gradual development is not as electrifying and showcase-worthy as are leaps of novelty, we can trace a large
%     amount elevation in knowledge that capitalises on this slow pace of progress.
%
% % TODO: margin note about the advancement in sensing gravitational waves <26-03-23, kunzaatko> 
%
% The minuteness of errors in our current observation has gotten to a scale that would be unimaginable for Kepler and
%     Brahe, but our desire for even greater precision prevails.
%
% \sidecite{kepler1609}
% % TODO: Add labels <20-05-23> 
%
% \section*{The History of Optical Instrumentation}%
% \labsec{optical_instrumentation}
% \addcontentsline{toc}{section}{Optical Instrumentation History}  
%
\pagelayout{margin}
